\documentclass[10pt,a4paper]{amsart}
\usepackage[margin=19mm]{geometry}
\usepackage{amsmath,amssymb,mathtools}
\usepackage[scr=rsfs,cal=euler]{mathalfa}
\usepackage{xfrac}
\usepackage[unicode,hidelinks,bookmarks=false]{hyperref}
\newcommand{\ie}{i.e.\ }
\newcommand{\cf}{cf.\ }
\newcommand{\resp}{respectively\ }
\newcommand{\Subsection}{\subsection}
\providecommand{\nobreakdash}{\nobreak\hbox{-}}
\setlength{\emergencystretch}{2em}
\multlinegap0pt
\usepackage[shortlabels]{enumitem}
\usepackage{amssymb}
\usepackage[normalem]{ulem}
\newtheorem{lemma}{Lemma}
\newtheorem{proposition}{Proposition}
\def\cal{\mathcal}
\def\ker{{\mbox{\rm{Ker} }}}
\def\m{{\mathfrak m}} %%% m
\title{Realization of spaces of commutative rings}
\author{Laura Cossu, Bruce Olberding}
\date{Source: arXiv:2409.18207v2 (CC BY 4.0)}
\begin{document}
\maketitle

\noindent\emph{Source and licence.} Realization of spaces of commutative rings. Version arXiv:2409.18207v2 (CC BY 4.0), distributed by arXiv under the Creative Commons Attribution 4.0 International licence, http://creativecommons.org/licenses/by/4.0/, as recorded in the arXiv licence metadata for this exact version and shown on the abstract page https://arxiv.org/abs/2409.18207v2. Full licence notice: \texttt{licenses/arxiv-2409.18207-CC-BY-4.0.txt}.\par\smallskip

\noindent\textit{Editorial scope.} Counted unit: the article's lemma patch (source0.tex lines 765--767). The article's complete original proof of it is delivered with it (source0.tex lines 768--800, one \texttt{proof} environment of 33 lines). No mathematics is written, repaired or re-derived: the statement and the proof are the article's own lines, in the article's own notation, and no retained line is substituted or re-flowed.  Retained uncounted as the source-local context the proof needs: lines 530--532; lines 710--719; lines 730--735.  Nothing was omitted from any retained span, so the package carries no omission record.  No retained line contains a citation, so the wrapper carries no bibliography and no bibliography entry.  No retained line contains a figure environment, an image-inclusion command or drawing code, so no image is delivered and the package declares no asset dependency.  Editorial formatting only, in addition: an AMS \texttt{amsart} preamble that restates the article's own declarations and macros used by the retained lines, namely the environments \texttt{lemma}, \texttt{proposition} on separate counters, together with the standard packages those lines need.  The retained environments print in this document as Lemma 1, Proposition 1 in order of appearance, whereas the article prints the same items with the numbering of its own full document.  The complete native source of the article as distributed in the arXiv e-print for this exact version is bundled unchanged as \texttt{sources/165890/source0.tex}.
\begin{lemma} \label{lem: bundle and stalks}
The image of a patch bundle $f:Y \rightarrow \Sigma(R)^{\rm zar}$  is the set of stalks of the patch presheaf ${\bf R}(f)$.    
\end{lemma} 

\begin{proposition} \label{prop: R presheaf}
Let $({\cal R},B)$ be a patch presheaf for the ring $R$. 

\begin{enumerate}[label=\textup{(\arabic{*})}, mode=unboxed]
\item\label{prop: R presheaf 1}
${\cal R}[B] $ is an $R_1$-algebra, where $1$ is the top element of $B$. 

\item\label{prop: R presheaf 2} $B$ is a Boolean algebra of faithful idempotents of ${\cal R}[B]$.
\end{enumerate}
\end{proposition}

\begin{lemma}\label{lem: R presheaf} Let $({\cal R}, B)$ be a patch presheaf for the ring $R$.  
\begin{enumerate}[label=\textup{(\roman{*})}, mode=unboxed]
    \item\label{lem: R presheaf i} If $a \in {\cal R}[B] $, then $a =  t_1e_1+ \cdots + t_ne_n,$ where $\{e_1,\ldots,e_n\}$ is an orthogonal subset of $B$, and $t_i \in R_{e_i}.$
    \item\label{lem: R presheaf ii} If $\{e_1,\ldots,e_n\}$ is an orthogonal subset of $B$ and $t_1,\ldots,t_n \in R$ are such that $t_1e_1+\cdots +t_ne_n  \in {\cal R}[B] $, then $t_i \in R_{e_i}$ for each $i$. 
\end{enumerate}
\end{lemma}

\begin{lemma} \label{lem: basic patch}
Let $A$ be the patch algebra of a patch bundle $f:Y \rightarrow X$ over $R$. For each point $y \in Y$, there is a homomorphism of $A$ onto the ring $f(y)$, whose kernel is an idempotent generated ideal of $A$.   
\end{lemma}

\begin{proof}
Let $({\cal R},B)$ be the patch presheaf that is the image of the patch bundle $f$ under the functor ${\bf R}$. Specifically, $B = {\rm Clop}(Y)$. Let ${y} \in Y$. 
We define a ring homomorphism $\psi_{y}:A \rightarrow f(y)$  in the following way. First, let $\m$ be the maximal ideal corresponding to $y$ under Stone duality, i.e., $$\m = \{U \in {\rm Clop}(Y):y \not\in U\}=\{e\in B : y\not \in e\}.$$ By Lemma~\ref{lem: R presheaf}\ref{lem: R presheaf i}, every element $a\in A$ admits a full orthogonal decomposition $a=\sum_i t_ie_i \in A$, where the $e_i$ form a full orthogonal set in $B$ and the coefficients $t_i \in R_{e_i}$. Since $\m$ is a maximal ideal of $B$ and this decomposition is full orthogonal, it follows that exactly one of the $e_i$ is not in $\m$. 
With this said, we define now 
\[\psi_{y}(a) = t_i,\text{ where }i\text{ is such that }e_i \not \in \m.\] By Lemma~\ref{lem: bundle and stalks}, $X$ is the set of stalks of $({\cal R},B)$, so 
$ f (y)={\mathcal{R}}_\m=\bigcup_{e\not \in \m} R_e$ and hence  $\psi_{y}(a) \in {\mathcal{R}}_{\m}$. 

To see that $\psi_y(a)$ is well defined, suppose $$a = \sum_i t_i e_i = \sum_j s_j f_j$$ are full orthogonal decompositions of $a$ with $t_i \in R_{e_i}$ and $s_j \in R_{f_j}$.  Without loss of generality, we may assume $e_1,f_1 \not \in \m$ and $e_i,f_j \in \m$ for all $i,j \geq 2$.   Then $$ae_1f_1 =t_1e_1f_1=s_1e_1f_1.$$  
If $e_1f_1 =0$, then at least one of $e_1,f_1$ is in $\m$, contrary to assumption. Thus $e_1f_1 \ne 0$ with $(t_1 -s_1)e_1f_1=0$. From the faithfulness of $e_1f_1$ (see Proposition~\ref{prop: R presheaf}\ref{prop: R presheaf 2}) conclude that $t_1 = s_1$, proving that $\psi_y$ is well defined. 

We claim next that $\psi_{y}$ is a ring homomorphism. 
Let $a,b \in A$ and let $\sum a_j f_j$ and $\sum b_k f'_k$ be full orthogonal decompositions of $a$ and $b$, respectively. By multiplying $a$ by $\sum f'_k=1$ and $b$ by $\sum f_j=1$, we may refine the full orthogonal decompositions of $a$ and $b$ to full orthogonal decompositions sharing the same idempotents, say,  $a =\sum_i t_i e_i$ and $b = \sum_i s_ie_i$, with $t_i,s_i \in R_{e_i}$. As noted above, we may assume without loss of generality that $e_1\not \in \m$ and $e_i \in \m$ for all $i \geq 2$.  By definition, \begin{center} $\psi_y(a+b) = t_1+s_1 = \psi_y(a)+\psi_y(b)$ and $\psi_y(ab) = t_1s_1 = \psi_y(a)\psi_y(b)$.
\end{center}
Also, $\psi_y(1) =1$,  and so $\psi_y$ is a ring homomorphism.



To show that $\psi_y$ is onto, pick an arbitrary element $r\in f(y)={\cal R}_\m=\bigcup_{e\not \in \m} R_e$ (note that in our notation $e$ is a clopen set of $Y$). Then $r\in R_{e_0}$ for some clopen $e_0$ not in $\m$ and $a=r e_0=\psi_y^{-1}(r)$.


We claim finally that  the kernel of $\psi_y$ is $\m A$, i.e., the ideal of $A$ generated by $\m$.
%\[\m A=\{e_1 a_1 + \dots + e_n a_n : e_i\in \m, a_i\in A, n\in \mathbb N\}.\]
It is easy to see that 
\begin{equation}\label{mA}
\m A = \bigcup_{e\in \m} eA.
\end{equation}
This follows from the fact that, if $e$ and $f$ are elements of $\m$ and $a, b\in A$, then $ea+fb=(e\vee f)(ea+fb)$, where $e\vee f \in \m$.
  

Let $a \in \ker \psi_y$, and write $a = \sum_i t_ie_i$, where the $e_i$ form a full orthogonal set of $B$ and $t_i \in R_{e_i}$ for every $i$. As noticed above, there is exactly one choice of $i$ for which $e_i \not \in \m$. For this choice of $i$, $t_i=0$, and so $a \in \m A$.  This proves $\ker \psi_y \subseteq \m A$. 

To prove the reverse inclusion, let $a \in \m A$. In light of \eqref{mA}, there is $e\in \m$ such that $a = e  \sum_i t_ie_i$, where the $e_i$ are pairwise orthogonal elements of $B$ and $t_i \in R_{e_i}$. Then, $a=\sum t_i f_i$, with $f_i=ee_i$ and $t_i\in R_{e_i}\subseteq R_{f_i}$. Since the $f_i$ are pairwise orthogonal and  all in $\m$, $\psi_y(a) =0$, proving that $\m A \subseteq \ker \psi_y$.
\end{proof} 

\end{document}
